\documentclass{beamer}
\usepackage{graphicx}
\usepackage{xcolor}
\usepackage{tikz}
\usepackage{amsmath}

\definecolor{purpleblue}{RGB}{80, 80, 180}
\definecolor{novelPurple}{RGB}{98, 54, 255}
\definecolor{softPurple}{RGB}{180, 170, 255}
\definecolor{softGray}{RGB}{235, 235, 245}
\definecolor{lightText}{RGB}{90, 90, 110}

\usetheme{Madrid}
\usecolortheme{dolphin}
\graphicspath{{static/course_materials/}{static/hard/}{./}}

\setbeamerfont{title}{size=\Large,series=\bfseries}
\setbeamerfont{frametitle}{size=\LARGE,series=\bfseries}
\setbeamercolor{frametitle}{fg=white, bg=novelPurple}
\setbeamercolor{title}{fg=white, bg=purpleblue}

\title[SAT Hard Question 5]{\textbf{SAT Hard Question\\ Set 5}}
\author{\textbf{Jack Zeng}}
\institute{\textcolor{white}{\textbf{Novel Prep}}}
\date{\today}

\begin{document}

\begin{frame}
    \titlepage
    \vfill
    \begin{flushright}
        \textcolor{purpleblue}{\Huge \textbf{Novel Prep}}
    \end{flushright}
\end{frame}

\begin{frame}
    \begin{tikzpicture}[remember picture, overlay]
        \shade[top color=softPurple, bottom color=softGray]
            (current page.north west) rectangle (current page.south east);
        \node[align=center] at (current page.center) {
            {\Huge\bfseries\textcolor{purpleblue}{Novel Prep}} \\[0.5cm]
            {\Large\itshape\textcolor{lightText}{Phase 2 · Hard Question Practice}}
        };
    \end{tikzpicture}
\end{frame}

\begin{frame}{Overview}
    \tableofcontents
\end{frame}

\begin{frame}
\section{Hard Question Set 5}
    \begin{tikzpicture}[remember picture, overlay]
        \shade[top color=softPurple, bottom color=softGray]
            (current page.north west) rectangle (current page.south east);
        \node[align=center] at (current page.center) {
            {\LARGE\bfseries\textcolor{purpleblue}{Hard Question Set 5}} \\[0.5cm]
            {\Large\itshape\textcolor{lightText}{15 SAT-style challenge problems}}
        };
    \end{tikzpicture}
\end{frame}

% --- Question 1 ---
\begin{frame}{Question 1}
\small
The given expression can be rewritten as \(\dfrac{a}{4}x^2 + \dfrac{b}{4}x + \dfrac{c}{4}\), where \(a, b,\) and \(c\) are constants. What is the value of \(a + b + c\)?

Given the expression
\[
-16(5x - 3)^2 + 4(5x - 2)^2
\]
\vspace{0.35cm}
\begin{enumerate}
    \item[A.] -112
    \item[B.] -56
    \item[C.] -28
    \item[D.] -7
\end{enumerate}
\end{frame}
\begin{frame}{Answer 1}
\small
\textbf{Correct Answer:} A
\end{frame}

% --- Question 2 ---
\begin{frame}{Question 2}
\small
In the given function, \(b\) is a positive constant. For every increase in the value of \(x\) by 1, the value of \(f(x)\) increases by \(c\%\), where \(0 < c < 100\).

Which expression gives the value of \(c\) in terms of \(b\)?
\[
f(x) = 66(b)^x
\]
\vspace{0.35cm}
\begin{enumerate}
    \item[A.] \(1 + \dfrac{b}{100}\)
    \item[B.] \(b + 100\)
    \item[C.] \(100(b + 1)\)
    \item[D.] \(100(b - 1)\)
\end{enumerate}
\end{frame}
\begin{frame}{Answer 2}
\small
\textbf{Correct Answer:} D
\end{frame}

% --- Question 3 ---
\begin{frame}{Question 3}
\small
If \(\sin(x^\circ) = a\), which of the following must be true for all values of \(x\)?
\vspace{0.35cm}
\begin{enumerate}
    \item[A.] \(\cos(x^\circ) = a\)
    \item[B.] \(\sin(90^\circ - x^\circ) = a\)
    \item[C.] \(\sin(x^2)^\circ = a^2\)
    \item[D.] \(\cos(90^\circ - x^\circ) = a\)
\end{enumerate}
\end{frame}
\begin{frame}{Answer 3}
\small
\textbf{Correct Answer:} D
\end{frame}

% --- Question 4 ---
\begin{frame}{Question 4}
\small
The function \(g\) is a quadratic function in the \(xy\)-plane. The graph of \(y = g(x)\) has a vertex at \((-1, -4)\) and passes through the points \((-2, -43)\) and \((1, -160)\).

What is the value of \(g(0) - g(2)\)?
\vspace{0.35cm}
\begin{enumerate}
    \item[A.] 312
    \item[B.] 355
    \item[C.] -117
    \item[D.] -203
\end{enumerate}
\end{frame}
\begin{frame}{Answer 4}
\small
\textbf{Correct Answer:} A
\end{frame}

% --- Question 5 ---
\begin{frame}{Question 5}
\small
The mass of object A is 324\% of the mass of object B and the mass of object A is 7.2\% of the mass of object C.

If the mass of object C is \(p\%\) of the mass of object B, what is the value of \(p\%\)?
\vspace{0.35cm}
\begin{enumerate}
    \item[A.] 4.5\%
    \item[B.] 4.5
    \item[C.] 45
    \item[D.] 450
\end{enumerate}
\end{frame}
\begin{frame}{Answer 5}
\small
\textbf{Correct Answer:} C
\end{frame}

% --- Question 6 ---
\begin{frame}{Question 6}
\small
In the given equation, \(c\) is a constant. If the equation has exactly one solution, what is the value of \(c\)?

\[
\frac{1}{cx} = \frac{x}{76} + \frac{1}{c}
\]
\end{frame}
\begin{frame}{Answer 6}
\small
\textbf{Final Answer:} $\boxed{-19}$
\end{frame}

% --- Question 7 ---
\begin{frame}{Question 7}
\small
A parallel electric circuit contains four resistors with resistances of \(p\) ohms, \(q\) ohms, \(s\) ohms, and 13 ohms.

The given equation relates the resistance \(R\), in ohms, of the circuit to the resistances of the four individual resistors it contains.

Which equation correctly expresses \(R\) in terms of \(p\), \(q\), and \(s\)?
\[
\frac{1}{R} = \frac{1}{p} + \frac{1}{q} + \frac{1}{s} + \frac{1}{13}
\]
\vspace{0.35cm}
\begin{enumerate}
    \item[A.] \(R = \dfrac{13pqs}{p + q + s + 13}\)
    \item[B.] \(R = \dfrac{13pqs}{13qs + 13ps + 13pq + pqs}\)
    \item[C.] \(R = \dfrac{13qs + 13ps + 13pq + pqs}{13pqs}\)
    \item[D.] \(R = p + q + s + 13\)
\end{enumerate}
\end{frame}
\begin{frame}{Answer 7}
\small
\textbf{Correct Answer:} B
\end{frame}

% --- Question 8 ---
\begin{frame}{Question 8}
\small
In the given equation, \(b\) is an integer. The equation has no real solutions.

What is the least possible value of \(b\)?
\[
-x^2 + bx - 169 = 0
\]
\vspace{0.35cm}
\begin{enumerate}
    \item[A.] 26
    \item[B.] 25
    \item[C.] -26
    \item[D.] -25
\end{enumerate}
\end{frame}
\begin{frame}{Answer 8}
\small
\textbf{Correct Answer:} D
\end{frame}

% --- Question 9 ---
\begin{frame}{Question 9}
\small
The exponential function \(f\) is defined by \(f(x) = ab^x\), where \(a\) and \(b\) are positive constants.

If \(f(n - 1) = f(n) + \dfrac{82}{100}f(n - 1)\), where \(n\) is a constant, what is the value of \(b\)?
\vspace{0.35cm}
\begin{enumerate}
    \item[A.] 1
    \item[B.] 0.82
    \item[C.] 0.18
    \item[D.] 10
\end{enumerate}
\end{frame}
\begin{frame}{Answer 9}
\small
\textbf{Correct Answer:} C
\end{frame}

% --- Question 10 ---
\begin{frame}{Question 10}
\small
A rectangle is inscribed in a circle such that the length of the diagonal of the rectangle is twice the length of its shortest side.

The circumference of the circle is \(114\pi\) units. What is the area, in square units, of the rectangle?
\vspace{0.35cm}
\begin{enumerate}
    \item[A.] \(57\sqrt{2}\)
    \item[B.] \(57\sqrt{3}\)
    \item[C.] \(3,\!249\sqrt{2}\)
    \item[D.] \(3,\!249\sqrt{3}\)
\end{enumerate}
\end{frame}
\begin{frame}{Answer 10}
\small
\textbf{Correct Answer:} D
\end{frame}

% --- Question 11 ---
\begin{frame}{Question 11}
\small
In triangle \(PQR\), the measure of angle \(P\) is \((3x + 5)^\circ\), the measure of angle \(Q\) is \((2x + 9)^\circ\), and the measure of angle \(R\) is \((4y + 5)^\circ\).

If side \(QR\) is extended through point \(R\) to point \(S\), and the measure of angle \(PRS\) is \((x + y)^\circ\), what is the value of \(x + y\)?
\vspace{0.35cm}
\begin{enumerate}
    \item[A.] 39
    \item[B.] 29
    \item[C.] 141
    \item[D.] 146
\end{enumerate}
\end{frame}
\begin{frame}{Answer 11}
\small
\textbf{Correct Answer:} A
\end{frame}

% --- Question 12 ---
\begin{frame}{Question 12}
\small
The manager of a gym selected a sample of 139 members at random to estimate the percentage of the gym's members that would continue to pay for a membership if the price increased.

From the survey, the manager estimates that 82\% of the gym's members would continue to pay for a membership if the price increased, with an associated margin of error of 6.39\%.

If the survey is repeated with a random sample of 278 members and the results are calculated in the same way, which of the following will be the most likely effect of using the larger random sample compared to the smaller random sample?
\vspace{0.35cm}
\begin{enumerate}
    \item[A.] The margin of error will be lower.
    \item[B.] The margin of error will be higher.
    \item[C.] The estimate of the percentage of the gym's members that would continue to pay for a membership if the price increased will be lower.
    \item[D.] The estimate of the percentage of the gym's members that would continue to pay for a membership if the price increased will be higher.
\end{enumerate}
\end{frame}
\begin{frame}{Answer 12}
\small
\textbf{Correct Answer:} A
\end{frame}

% --- Question 13 ---
\begin{frame}{Question 13}
\small
Given the expression \(\sqrt[5]{119n \left(\sqrt[6]{119n}\right)^2}\),

for what value of \(x\) is the expression equivalent to \((119n)^{30x}\), where \(n > 1\)?
\vspace{0.35cm}
\begin{enumerate}
    \item[A.] 225
    \item[B.] \(\dfrac{8}{225}\)
    \item[C.] \(\dfrac{4}{225}\)
    \item[D.] \(\dfrac{2}{225}\)
\end{enumerate}
\end{frame}
\begin{frame}{Answer 13}
\small
\textbf{Correct Answer:} D
\end{frame}

% --- Question 14 ---
\begin{frame}{Question 14}
\small
The mass of object A is 483\% of the mass of object B, and the mass of object A is 0.069\% of the mass of object C.

If the mass of object C is \(p\%\) of the mass of object B, what is the value of \(\dfrac{p}{1,\!000}\)?
\vspace{0.35cm}
\begin{enumerate}
    \item[A.] 0.7\%
    \item[B.] 70
    \item[C.] 70\%
    \item[D.] 700
\end{enumerate}
\end{frame}
\begin{frame}{Answer 14}
\small
\textbf{Correct Answer:} D
\end{frame}

% --- Question 15 ---
\begin{frame}{Question 15}
\small
For triangle \(PQR\) and triangle \(STU\), angles \(P\) and \(S\) each measure \(30^\circ\), \(QR = 16\), and \(TU = 12\).

Which additional piece(s) of information is(are) sufficient to prove that triangle \(PQR\) is similar to triangle \(STU\)?

\begin{itemize}
\item[I] \(\angle Q = 40^\circ\) and \(\angle U = 110^\circ\)
\item[II] \(\dfrac{PR}{SU} = \dfrac{4}{3}\)
\end{itemize}
\vspace{0.35cm}
\begin{enumerate}
    \item[A.] I is sufficient, but II is not.
    \item[B.] II is sufficient, but I is not.
    \item[C.] I is sufficient, and II is sufficient.
    \item[D.] Neither I nor II is sufficient.
\end{enumerate}
\end{frame}
\begin{frame}{Answer 15}
\small
\textbf{Correct Answer:} C
\end{frame}

\begin{frame}{}
    \begin{tikzpicture}[remember picture, overlay]
        \shade[top color=novelPurple!40, bottom color=white]
            (current page.north west) rectangle (current page.south east);
        \fill[white, opacity=0.15] (current page.center) circle (5cm);
        \foreach \x in {1,2,3,4,5} {
            \fill[novelPurple, opacity=0.12] (rand*10-5, rand*7-3.5) circle (0.1+\x*0.05);
        }
    \end{tikzpicture}
    \centering
    \vspace{5em}
    {\Huge \textbf{\textcolor{novelPurple}{Thank You!}}} \\[1em]
    {\large \textcolor{gray}{We appreciate your time and focus.}} \\[3em]
    {\LARGE \textbf{\textcolor{novelPurple}{\textsf{Novel Prep}}}}
    \vfill
\end{frame}

\end{document}
